\documentclass{proc}

\usepackage[utf8]{inputenc}
\usepackage{amsmath, amssymb, amsfonts}
\usepackage{booktabs}
\usepackage{array}
\usepackage{microtype}
\usepackage{cite}

\title{Precision Gauge Theory Simulations and Non-Perturbative Parton Structure at the Meridian High-Energy Facility}

\author{Elena Vance \and Marcus Sterling \\
Department of High Energy Physics, Nexa Institute of Technology \\
\and
Julian Thorne \and Clara Rey \\
Center for Computational Subatomic Physics, Vertex Research Laboratory \\
\and
David K. Chen \\
Synapse Quantum Science Initiative, Meridian Science Park}

\begin{document}

\maketitle

\begin{abstract}
We present high-precision non-perturbative calculations of hadronic matrix elements and parton distribution functions derived from lattice gauge theory simulations at physical quark masses. Using gauge field configurations generated on the Vertex Parallel Supercomputing Cluster with a tree-level improved Symanzik gauge action and $N_f = 2+1+1$ flavors of highly improved staggered quarks (HISQ), we compute the pion and nucleon vector and axial-vector form factors across momentum transfers up to $Q^2 = 3.0~\text{GeV}^2$. Continuum extrapolations are performed across three lattice spacings ($a \approx 0.06, 0.09, 0.12~\text{fm}$) with physical spatial volumes satisfying $m_\pi L > 4.0$. Hadronic matrix elements are renormalized non-perturbatively via the regularized $\text{RI}^\prime\text{-MOM}$ scheme and converted to the $\overline{\text{MS}}$ scheme at $2~\text{GeV}$. Our results for the nucleon axial charge $g_A = 1.268(14)$ and the pion decay constant $f_\pi = 130.4(9)~\text{MeV}$ demonstrate sub-percent control over systematic uncertainties, providing key theoretical inputs for ongoing and future scattering experiments at the Meridian Collider Facility.
\end{abstract}

\section{Introduction}
Understanding the non-perturbative structure of hadrons from first principles in Quantum Chromodynamics (QCD) remains one of the primary goals of modern subatomic physics. While high-energy hard scattering processes are well described by perturbative QCD expansions in powers of the strong coupling constant $\alpha_s(Q^2)$, low-energy observables and bound-state properties necessitate non-perturbative techniques. Lattice QCD provides a mathematically rigorous, non-perturbative regularization of the gauge theory formulated on a four-dimensional Euclidean spacetime grid.

In recent years, advances in algorithmic design---such as adaptive algebraic multigrid Dirac solvers and Fourier-accelerated Hybrid Monte Carlo (HMC) algorithms---combined with high-performance computing architectures have enabled lattice simulations to operate directly at physical light-quark masses. This eliminates the reliance on chiral extrapolations from unphysically heavy pion masses, which historically served as a major source of systematic error.

This report summarizes recent progress achieved by the Nexa-Vertex-Synapse Lattice Collaboration in evaluating hadron matrix elements relevant for lepton-hadron scattering experiments conducted at the Meridian High-Energy Collider. Special emphasis is placed on controlling continuum, finite-volume, and excited-state contamination.

\section{Theoretical Framework}
The Euclidean action for lattice gauge theory consists of the pure gauge action $S_G[U]$ and the fermion action $S_F[\psi, \bar{\psi}, U]$:
\begin{equation}
S[U, \psi, \bar{\psi}] = S_G[U] + S_F[\psi, \bar{\psi}, U].
\end{equation}

For the gauge sector, we employ the tree-level improved Symanzik action defined on the hypercubic lattice $\Lambda = \{ x = (x_1, x_2, x_3, x_4) \mid x_\mu / a \in \mathbb{Z} \}$:
\begin{align}
S_G[U] = \frac{\beta}{3} \sum_{x, \mu < \nu} &\Big( c_0 \text{Re Tr}\left[ 1 - P_{\mu\nu}(x) \right] \nonumber \\
&+ c_1 \text{Re Tr}\left[ 1 - R_{\mu\nu}(x) \right] \Big),
\end{align}
where $P_{\mu\nu}(x)$ denotes the standard $1 \times 1$ plaquette loop, $R_{\mu\nu}(x)$ represents the rectangular $1 \times 2$ loop, $\beta = 6/g^2$ is the lattice gauge coupling, and the normalization condition requires $c_0 + 8c_1 = 1$ with $c_1 = -1/12$.

The fermion action uses the $N_f = 2+1+1$ Highly Improved Staggered Quark (HISQ) formulation for sea quarks and the Wilson-clover action for valence u, d, and s quarks. The Wilson-clover Dirac operator $D_{\text{clover}}$ is given by:
\begin{align}
D_{\text{clover}}(x, y) &= \left( m_0 + \frac{4}{a} \right) \delta_{x,y} \nonumber \\
&-\frac{1}{2a} \sum_{\mu=1}^4 \Big[ (1 - \gamma_\mu) U_\mu(x) \delta_{x+a\hat{\mu}, y} \nonumber \\
&\quad + (1 + \gamma_\mu) U_\mu^\dagger(y) \delta_{x-a\hat{\mu}, y} \Big] \nonumber \\
&+ c_{\text{SW}} \frac{a}{4} \sum_{\mu,\nu} \sigma_{\mu\nu} F_{\mu\nu}(x) \delta_{x,y},
\end{align}
where $c_{\text{SW}}$ is the Sheikholeslami-Wohlert coefficient tuned non-perturbatively to eliminate $\mathcal{O}(a)$ discretization errors, and $F_{\mu\nu}(x)$ is the lattice field strength tensor constructed from clover-shaped four-plaquette loops.

\section{Computational Methodology}
Gauge field ensembles were generated using the Hasenbusch-preconditioned Hybrid Monte Carlo algorithm integrated with an adaptive leapfrog integrator. The numerical inversions of the Dirac operator were accelerated using an algebraic multigrid preconditioner implemented on the Apex-64 GPU cluster hosted at the Vertex Research Laboratory.

Table~\ref{tab:ensembles} outlines the key parameters of the five gauge configurations used in this analysis. Spatial volumes $L^3 \times T$ were chosen such that $m_\pi L \ge 4.0$ across all ensembles to suppress finite-volume effects to below $0.2\%$.

\begin{table}[t]
\caption{Parameters of the $N_f = 2+1+1$ lattice gauge ensembles generated by the Nexa-Vertex collaboration. Listed are the lattice coupling $\beta$, volume $(L/a)^3 \times (T/a)$, lattice spacing $a$, pion mass $m_\pi$, and total HMC trajectories.}
\label{tab:ensembles}
\vspace{4pt}
\centering
\small
\setlength{\tabcolsep}{3pt}
\begin{tabular}{lccccc}
\toprule
Ensemble & $\beta$ & Volume & $a$ (fm) & $m_\pi$ (MeV) & Traj. \\
\midrule
NV-A06   & 6.30 & $64^3 \times 128$ & 0.058 & 134.2 & 8,400 \\{}
NV-A09a  & 6.00 & $48^3 \times 96$  & 0.088 & 135.1 & 10,200 \\{}
NV-A09b  & 6.00 & $32^3 \times 64$  & 0.088 & 220.4 & 12,000 \\{}
NV-A12a  & 5.80 & $32^3 \times 64$  & 0.121 & 136.0 & 9,600 \\{}
NV-A12b  & 5.80 & $24^3 \times 48$  & 0.121 & 310.8 & 14,500 \\
\bottomrule
\end{tabular}
\end{table}

Correlator measurements were carried out using momentum-smeared point and gaussian sources. Ground-state isolation was achieved by performing multi-state exponential fits to three-point correlation functions over multiple source-sink separations $\tau \in [0.8, 1.6]~\text{fm}$.

\section{Results and Error Budget}
Hadronic matrix elements $\langle H(p') | \mathcal{O}_{\mu} | H(p) \rangle$ were extracted from ratios of three-point to two-point correlation functions:
\begin{equation}
R_\mathcal{O}(t, \tau) = \frac{C^{(3)}_\mathcal{O}(t, \tau; \vec{p}', \vec{p})}{C^{(2)}(\tau; \vec{p}')} \left[ \mathcal{K}(t, \tau; \vec{p}', \vec{p}) \right]^{1/2},
\end{equation}
where the kinematic ratio factor $\mathcal{K}$ is defined as:
\begin{equation}
\mathcal{K} = \frac{C^{(2)}(t; \vec{p}') C^{(2)}(\tau; \vec{p}') C^{(2)}(\tau - t; \vec{p})}{C^{(2)}(t; \vec{p}) C^{(2)}(\tau; \vec{p}) C^{(2)}(\tau - t; \vec{p}')}.
\end{equation}

For the nucleon axial charge $g_A$, the asymptotic plateau value as $\tau \to \infty$ yields the bare matrix element. Renormalization factors $Z_A$ were determined non-perturbatively using the $\text{RI}^\prime\text{-MOM}$ scheme on the lattice and matched to $\overline{\text{MS}}$ at $\mu = 2.0~\text{GeV}$.

\begin{table}[t]
\caption{Systematic error budget for the nucleon axial charge $g_A$ and pion decay constant $f_\pi$ evaluated in this work.}
\label{tab:systematics}
\vspace{4pt}
\centering
\small
\setlength{\tabcolsep}{3.5pt}
\begin{tabular}{lcc}
\toprule
Source of Uncertainty & $\Delta g_A / g_A$ (\%) & $\Delta f_\pi / f_\pi$ (\%) \\
\midrule
Statistical & 0.72 & 0.45 \\{}
Continuum Limit ($a \to 0$) & 0.55 & 0.38 \\{}
Chiral Interpolation & 0.31 & 0.20 \\{}
Finite Volume ($m_\pi L$) & 0.18 & 0.12 \\{}
Renormalization ($Z_A, Z_V$) & 0.40 & 0.30 \\{}
Excited-State Effects & 0.50 & 0.25 \\
\midrule
Total Uncertainty & 1.10 & 0.69 \\
\bottomrule
\end{tabular}
\end{table}

Extrapolating to the continuum limit via a joint fit of the form:
\begin{align}
f(a, m_\pi) &= f_0 + c_a \left(\frac{a}{a_0}\right)^2 + c_\pi \left(\frac{m_\pi}{m_{\pi,\text{phys}}}\right)^2 \nonumber \\
&\quad + c_L \frac{e^{-m_\pi L}}{(m_\pi L)^{3/2}},
\end{align}
we obtain our final physical predictions:
\begin{align}
g_A &= 1.268 \pm 0.009_{\text{stat}} \pm 0.011_{\text{sys}}, \\
f_\pi &= 130.4 \pm 0.6_{\text{stat}} \pm 0.7_{\text{sys}}~\text{MeV}.
\end{align}

These results are in excellent agreement with experimental world averages while providing unprecedented precision on non-perturbative parameters required for Meridian Collider phenomenology.

\section{Summary and Outlook}
We have presented a non-perturbative calculation of hadron structure observables using high-statistics lattice QCD simulations on $N_f = 2+1+1$ HISQ ensembles. Through meticulous control of excited states and continuum extrapolations, sub-percent precision has been attained for key decay constants and form factors.

Future work will expand these calculations to flavor-singlet parton distributions and transverse-momentum-dependent (TMD) structure functions utilizing next-generation exascale compute allocations at the Apex Supercomputing Center.

\section*{Acknowledgments}
This work was supported by the Meridian Research Foundation under Grant No.~MRF-2024-88402, the Nexa Institute Computing Initiative, and Vertex Supercomputing Grant VS-4091. Computations were performed on the Apex-64 cluster at Vertex Research Laboratory and the Synapse HPC facility.

\begin{thebibliography}{99}
\bibitem{Vance2024}
E. Vance, M. Sterling, and J. Thorne, \textit{Non-perturbative lattice QCD at physical mass}, Phys. Rev. D \textbf{109}, 054501 (2024).

\bibitem{Sterling2023}
M. Sterling et al. (Nexa Collaboration), \textit{Hadron form factors from HISQ configurations}, PoS \textbf{LATTICE2023}, 142 (2023).

\bibitem{Chen2025}
D. K. Chen and C. Rey, \textit{High-precision renormalization in RI-MOM scheme}, J. High Energy Phys. \textbf{04}, 088 (2025).

\bibitem{Thorne2022}
J. Thorne, \textit{Algorithmic developments for GPU lattice solvers}, Comput. Phys. Commun. \textbf{275}, 108310 (2022).

\bibitem{MeridianReport2025}
Meridian Collider Collaboration, \textit{Physics opportunities at the Meridian High-Energy Facility}, Report No. MCC-2025-01 (2025).

\bibitem{Symanzik1983}
K. Symanzik, \textit{Continuum limit and improved action in lattice field theory}, Nucl. Phys. B \textbf{226}, 187 (1983).

\bibitem{Follana2007}
E. Follana et al. (HPQCD Collaboration), \textit{Highly improved staggered quarks on the lattice}, Phys. Rev. D \textbf{75}, 054502 (2007).
\end{thebibliography}

\end{document}
